\chapter{Mạng nơ-ron: từ perceptron đến lan truyền ngược}

\section{Perceptron và giới hạn của tuyến tính}

\emph{Perceptron} có thể được xem là mô hình phân loại tuyến tính có kích hoạt ngưỡng:
\[
\hat{y} = \sgn(\bm{w}^\top \bm{x}+b).
\]
Nếu dữ liệu tách tuyến tính, thuật toán perceptron có thể tìm một siêu phẳng tách. Nhưng khi quan hệ giữa đầu vào và đầu ra không tuyến tính, perceptron đơn lẻ thất bại. Bài toán XOR là ví dụ kinh điển.

Điều này dẫn đến câu hỏi tự nhiên: nếu một lớp tuyến tính không đủ, ta xếp nhiều lớp lại có đủ không? Câu trả lời là có, miễn là giữa các lớp ta đặt một ánh xạ phi tuyến. Chính từ đây, \emph{multi-layer perceptron} (mạng nhiều lớp) ra đời.

\section{Một lớp ẩn đã làm thay đổi tất cả}

Xét một mạng một lớp ẩn:
\[
\bm{h} = \phi(W_1\bm{x}+\bm{b}_1),\qquad
\hat{\bm{y}} = W_2\bm{h}+\bm{b}_2.
\]
Nếu \(\phi\) là affine, toàn bộ mạng vẫn chỉ là affine. Vì vậy, phi tuyến là bắt buộc. Các hàm kích hoạt phổ biến gồm \emph{sigmoid}, \(\tanh\), và \emph{ReLU} (hàm tuyến tính chỉnh lưu). Mỗi hàm có ưu và nhược điểm. \emph{Sigmoid} và \(\tanh\) trơn và dễ liên hệ tới xác suất/chuẩn hóa, nhưng dễ bão hòa. \emph{ReLU} đơn giản, gradient ổn ở miền dương, nhưng có thể ``chết'' ở miền âm.

\subsection{Định lý xấp xỉ phổ quát}

Một kết quả có sức cổ vũ lớn cho mạng nơ-ron là \emph{universal approximation theorem} (định lý xấp xỉ phổ quát): với số nơ-ron đủ lớn và phi tuyến phù hợp, mạng một lớp ẩn có thể xấp xỉ bất kỳ hàm liên tục nào trên tập compact đến độ chính xác tùy ý. Điều này không nói rằng học là dễ, không nói rằng cần ít nơ-ron, và không nói rằng tổng quát hóa sẽ tốt. Nó chỉ nói lớp hàm của MLP là rất phong phú.

Vì vậy, sức mạnh của mạng nơ-ron không đến từ một công thức thần bí, mà đến từ việc ta được quyền xây dựng các biểu diễn phi tuyến rất giàu.

\section{Hàm kích hoạt và hình học của gradient}

Hãy nhìn vào đạo hàm:
\[
\sigmoid'(z)=\sigmoid(z)(1-\sigmoid(z)),
\qquad
\tanhf'(z)=1-\tanh^2(z),
\qquad
\relu'(z)=\mathbf{1}_{z>0}.
\]
Với \emph{sigmoid} và \(\tanh\), khi \(|z|\) lớn thì đạo hàm gần 0. Trong mạng nhiều lớp, nhân lặp nhiều đạo hàm nhỏ dễ dẫn đến \emph{vanishing gradient} (triệt tiêu gradient). \emph{ReLU} giảm bớt vấn đề này trên miền dương, nhưng đổi lại không khả vi tại 0 và có miền âm vô hoạt.

Lựa chọn hàm kích hoạt vì vậy không phải vấn đề thẩm mỹ. Nó thay đổi động học học. Nhiều quy tắc thực hành trong học sâu thực ra là cách khống chế động học gradient.

\section{Suy dẫn lan truyền ngược từ quy tắc dây chuyền}

Ta xét mạng một lớp ẩn cho bài toán hồi quy, với hàm mất mát trên một mẫu:
\[
L = \frac{1}{2}\|\hat{\bm{y}}-\bm{y}\|_2^2,
\qquad
\hat{\bm{y}} = W_2\bm{h}+\bm{b}_2,
\qquad
\bm{h} = \phi(\bm{z}_1),
\qquad
\bm{z}_1 = W_1\bm{x}+\bm{b}_1.
\]
Đặt
\[
\bm{\delta}_2 = \frac{\partial L}{\partial \bm{z}_2}
=
\hat{\bm{y}}-\bm{y},
\]
vì \(\bm{z}_2=\hat{\bm{y}}\) trong lớp cuối affine. Khi đó
\[
\frac{\partial L}{\partial W_2} = \bm{\delta}_2 \bm{h}^\top,
\qquad
\frac{\partial L}{\partial \bm{b}_2}=\bm{\delta}_2.
\]
Để lan truyền ngược về lớp ẩn, ta dùng quy tắc dây chuyền:
\[
\bm{\delta}_1
=
\frac{\partial L}{\partial \bm{z}_1}
=
\left(W_2^\top \bm{\delta}_2\right)\odot \phi'(\bm{z}_1).
\]
Do đó
\[
\frac{\partial L}{\partial W_1}=\bm{\delta}_1 \bm{x}^\top,
\qquad
\frac{\partial L}{\partial \bm{b}_1}=\bm{\delta}_1.
\]

\begin{remark}
Lan truyền ngược không phải một thuật toán ``khác'' hạ dốc theo gradient. Nó là kỹ thuật tính gradient hiệu quả để sau đó dùng trong hạ dốc theo gradient hay bất kỳ bộ tối ưu nào khác.
\end{remark}

\section{Ma trận hóa trên \emph{mini-batch}}

Trong thực hành, ta huấn luyện trên \emph{mini-batch} (lô nhỏ) kích thước \(B\). Nếu đặt \(X\in\R^{d\times B}\) là ma trận gồm các vectơ đầu vào theo cột, thì
\[
Z_1=W_1X+\bm{b}_1 \mathbf{1}^\top,
\qquad
H=\phi(Z_1),
\qquad
\hat{Y}=W_2H+\bm{b}_2\mathbf{1}^\top.
\]
Đạo hàm trên cả lô là tổng hoặc trung bình đạo hàm từng mẫu. Thao tác ma trận hóa này không chỉ nhanh hơn về tính toán, mà còn bộc lộ rõ hơn cấu trúc tuyến tính của phép lan truyền tiến và lan truyền ngược.

\section{Từ MLP sang học biểu diễn}

Mỗi lớp ẩn có thể được xem là một phép biến đổi biểu diễn. Đầu vào ban đầu có thể là các đặc trưng thô, nhưng sau nhiều lớp, mạng tự tạo ra các đặc trưng trung gian hữu ích cho hàm mất mát cuối. Chính vì vậy, học sâu được mô tả là \emph{representation learning} (học biểu diễn).

Hãy hình dung bài toán phân loại ảnh chữ số viết tay. Một mô hình cổ điển cần ta thiết kế đặc trưng: nét ngang, nét dọc, điểm giao, v.v. MLP sâu trên dữ liệu ảnh có thể học những biến đổi trung gian đó tự động. Trong dự báo chuỗi thời gian, \emph{LSTM} sẽ học biểu diễn trạng thái tạm thời của quá trình sinh dữ liệu.

\section{Điều chuẩn trong mạng nơ-ron}

Vì mạng nơ-ron có sức biểu diễn lớn, quá khớp là nguy cơ thường trực. Các cơ chế điều chuẩn phổ biến gồm:

\begin{enumerate}
\item phạt chuẩn của hệ số;
\item \emph{early stopping} (dừng sớm);
\item \emph{dropout};
\item \emph{data augmentation} (tăng cường dữ liệu);
\item chia sẻ tham số và ràng buộc kiến trúc.
\end{enumerate}

\emph{Dropout} có thể được xem như một hình thức trung bình mô hình xấp xỉ: trong mỗi bước huấn luyện, một tập con nơ-ron bị tắt ngẫu nhiên. Điều này ép mạng không phụ thuộc quá mạnh vào bất kỳ đường dẫn nào.

\emph{Early stopping} là điều chuẩn đặc biệt quan trọng. Nhưng nó chỉ đúng nghĩa khi được ra quyết định trên tập thẩm định, không phải tập kiểm tra. Đây là một điểm nhỏ trong code nhưng lớn trong logic thống kê.

\section{Triệt tiêu và bùng nổ gradient}

Trong mạng sâu, gradient tại lớp đầu có thể chứa tích của nhiều Jacobian:
\[
\frac{\partial L}{\partial \bm{h}^{(1)}}
=
\left(\prod_{\ell=2}^{L} J_\ell^\top\right)
\frac{\partial L}{\partial \bm{h}^{(L)}}.
\]
Nếu các Jacobian này có chuẩn phổ biến nhỏ hơn 1, gradient giảm dần về 0. Nếu lớn hơn 1, nó có thể nổ. Vấn đề này biểu lộ rõ trong RNN, nơi cùng một ma trận được nhân lặp theo thời gian.

Nhiều kỹ thuật sau này thực ra là phản hồi với vấn đề này: khởi tạo tốt, \emph{ReLU}, chuẩn hóa theo lô, kết nối phần dư, \emph{LSTM gates} (các cổng của LSTM), và \emph{gradient clipping} (cắt gradient).

\section{\emph{Softmax} và entropy chéo cho phân loại nhiều lớp}

Nếu đầu ra có \(K\) lớp, ta đặt \emph{logit} \(\bm{z}\in\R^K\) và
\[
\hat{p}_k = \frac{e^{z_k}}{\sum_{j=1}^K e^{z_j}}.
\]
Với nhãn \emph{one-hot}, \emph{cross-entropy} là
\[
L = -\sum_{k=1}^K y_k \log \hat{p}_k.
\]
Một công thức quan trọng là
\[
\frac{\partial L}{\partial z_k} = \hat{p}_k-y_k.
\]
Công thức này làm cho lan truyền ngược của lớp cuối trở nên gọn gàng và là một lý do thực dụng khiến cặp \emph{softmax + cross-entropy} rất phổ biến.

\section{Tối ưu trong mạng nơ-ron}

Sau khi có gradient, ta cần một bộ tối ưu. Hạ dốc theo gradient toàn bộ dữ liệu quá tốn kém; do đó ta dùng \emph{stochastic} hay \emph{mini-batch gradient descent}. \emph{Momentum} thêm trí nhớ quán tính. \emph{Adam} chuẩn hóa gradient theo ước lượng moment bậc nhất và bậc hai. Các bộ tối ưu hiện đại không thay thế nhu cầu hiểu hàm mất mát và dữ liệu; chúng chỉ là công cụ lái bài toán tối ưu.

Một điểm cần ghi nhớ: bộ tối ưu không thể cứu một giao thức đánh giá sai. Nếu tập huấn luyện/tập thẩm định/tập kiểm tra bị dùng lẫn, dù dùng Adam hay SGD cũng không làm kết quả đáng tin hơn.

\section{Từ mạng nơ-ron đến tư duy hồi quy theo thời gian}

MLP xử lý đầu vào có kích thước cố định và không thể hiện rõ cấu trúc thứ tự. Chuỗi thời gian, văn bản, âm thanh lại có bản chất tuần tự. Nếu ta cắt chuỗi thành vectơ và ném vào MLP, ta có thể mất thông tin thứ tự. Điều này dẫn đến \emph{recurrent neural network} (mạng nơ-ron hồi tiếp), nơi tham số được chia sẻ qua các bước thời gian:
\[
\bm{h}_t = \phi(W_x \bm{x}_t + W_h \bm{h}_{t-1} + \bm{b}).
\]
Nhưng ở đây, vấn đề triệt tiêu/bùng nổ gradient trở nên nghiêm trọng hơn vì nhân lặp cùng một ma trận trong nhiều bước. LSTM được sinh ra để giải quyết điều này. Chương về chuỗi thời gian sẽ quay lại kỹ hơn.

\section{Mạng nơ-ron và vấn đề diễn giải}

MLP mạnh về biểu diễn nhưng yếu về minh bạch. Ta có thể xem gradient, activation, \emph{feature importance} (mức quan trọng của đặc trưng) cục bộ, nhưng không dễ đọc một mạng như đọc một cây quyết định hay một bộ luật mờ. Vì vậy, các nỗ lực lai ghép fuzzy và mạng nơ-ron có sức hút tự nhiên. Tuy nhiên, bài học từ mạng nơ-ron là: chỉ cần một lớp kết hợp cuối cùng đủ phức tạp, khả năng diễn giải toàn cục đã giảm mạnh. Chúng ta sẽ phải nhớ điều này khi nhìn mô hình lai \emph{fuzzy + BiLSTM}.

\section{Kết}

Chương này cần được nhớ bằng một ý tưởng cốt lõi: mạng nơ-ron không phải hộp đen vì nó không có toán học, mà vì toán học của nó được tổ chức trong một không gian tham số lớn, phi tuyến, và khó đọc. Lan truyền ngược dựa trên quy tắc dây chuyền. Sức mạnh của nó đến từ học biểu diễn. Vấn đề của nó đến từ tối ưu, tổng quát hóa, và diễn giải. Toàn bộ ba vấn đề ấy sẽ trở lại nhiều lần trong phần học sâu và ANFIS.
